\section{走向大规模多模态 Transformer}

前面的模型分别处理空间基因、H\&E 条件或局部邻域。下一阶段要把这些模块放进同一患者表示，并允许模型在不同模态间翻译、消歧和模拟干预。本章比较单模态与多模态架构，随后解释一个多模态 counterfactual demo。

\subsection{统一数据与统一主干}

统一模型的诱惑是让所有数据都进入一个大 Transformer；真正难点却是 token 数量、缺失模式、分辨率和监督信号极不均衡。若设计不当，最便宜或 token 最多的模态会支配训练。下面先用单模态模型建立可诊断 baseline，再观察共享层是否真正带来跨模态增益，而不是只增加参数量。

\lecturefigure{slide-237-dataset-synthesis.jpg}{数据汇总：H\&E、蛋白、空间转录组与遗传信息}{00:50:19}

这张页与前面的四模态页相似，但角色不同：前者用于介绍数据，当前页用于决定模型输入。H\&E 和蛋白是稠密空间图像，空间转录组是稀疏位置上的高维计数，遗传信息更接近患者级条件。统一前必须明确各自的 tokenization 与 mask policy。

\lecturefigure{slide-238-unimodal-transformer.jpg}{单模态 Transformer：一个输入流对应一个预测头}{00:50:25}

单模态路线易训练、易诊断，也适合作为 baseline。它能回答“仅凭这一模态能学到什么”，帮助测量新增模态的边际价值。若多模态模型只比最佳单模态略好，可能说明数据对齐不足、融合未生效，或新增模态本身不含任务信息。

\lecturefigure{slide-240-multimodal-transformer.jpg}{多模态 Transformer：多个编码流汇入共享融合层和输出空间}{00:51:21}

图中多种颜色代表不同模态 token，中央 multimodal layers 负责交互。输出可重建任一模态，也可预测治疗相关目标。训练时要平衡不同损失，否则容易出现“重建 H\&E 很好、关键分子任务很差”的目标错配。

\begin{knowledgebox}{Masking policy 是隐含课程设计}
遮蔽哪些 token、遮蔽多少、是否整模态缺失，决定模型被迫学习什么。随机遮少量基因主要学习局部相关；整块遮掉一个模态才真正训练跨模态 translation；同时遮空间区域则更接近组织补全与远距离推理。
\end{knowledgebox}

\subsection{多模态反事实模拟}

统一模型的核心承诺是：改变一个模态或干预条件后，其他模态的预测也应一致更新。这比单独生成一张表达图更接近世界模型，但仍需要严格验证干预是否遵循真实机制。为了读懂接下来的 demo，需要固定患者与空间上下文，只比较被修改条件前后的联合输出，并把图像变化与可测实验终点对应起来。

\lecturefigure{slide-241-counterfactual-architecture.jpg}{多模态 counterfactual：改变输入模态并观察联合预测变化}{00:52:00}

左侧输入包含组织图像、空间测量和可能的治疗条件，中间是共享模型，右侧输出多模态患者表示。理想情况下，模型能回答“若改变该基因或治疗，蛋白、RNA 和组织状态将怎样共同变化”。模型若只做静态补全，则不具备这种一致的状态转移。

\lecturefigure{slide-247-counterfactual-results.jpg}{反事实 demo 的最终状态：输入改变后输出地图出现空间差异}{00:53:36}

读图先看哪些输入被固定，再看被修改模态和输出差异。红绿区域显示预测方向，但不直接代表疗效。可靠验证需要已知 perturbation 数据：把真实干预前信息输入模型，比较模型预测与干预后测量，而不是只判断图是否“像生物学”。

\begin{lstlisting}[language=Python,caption={实验优先级：把模型效应、不确定性与现实验证成本同时纳入}]
def rank_experiments(candidate_interventions, patient_context):
    ranked = []
    for intervention in candidate_interventions:
        effect, uncertainty = world_model.simulate(
            patient_context, intervention
        )
        novelty = compare_with_existing_evidence(intervention)
        cost = wet_lab_cost(intervention)
        score = expected_value(effect, uncertainty, novelty, cost)
        ranked.append((score, intervention))
    return sorted(ranked, reverse=True)
\end{lstlisting}

\begin{importantbox}{模拟器与实验选择器是两个系统}
世界模型估计“如果做 $a$ 会怎样”；scientist 或 agent 决定“哪个 $a$ 值得做”。即使模拟器完美，穷举所有组合也不现实。实验价值还取决于新颖性、可验证性、成本与安全。
\end{importantbox}

\subsection{本章小结}

大规模多模态 Transformer 需要协调异构 token、缺失模态、masking policy 和多目标损失。单模态模型是必要 baseline；统一模型的增量价值应体现在更好的 translation、disambiguation 与干预一致性上。Counterfactual demo 是候选机制生成器，不是临床疗效证明。

